\chapter{Visao geral da pesquisa}

\section{Problema e recorte}

A arquitetura Transformer sustenta grande parte dos sistemas modernos de
inteligencia artificial, mas suas projecoes lineares densas e o mecanismo de
self-attention concentram custo de processamento, memoria e movimentacao de
dados. Esta iniciacao cientifica investiga como configuracoes numericas e
estruturais alteram esse custo durante a inferencia.

O objeto experimental foi deliberadamente reduzido a duas operacoes centrais:
projecao linear densa e scaled dot-product attention. O codigo inclui um bloco
Transformer simplificado para validacao arquitetural, mas nao reivindica uma
arquitetura Transformer completa. Essa delimitacao permite isolar efeitos,
repetir os experimentos e relacionar cada resultado ao mecanismo que o produziu.

\section{Objetivo}

Analisar, em cenarios controlados de inferencia, os compromissos entre eficiencia
computacional e preservacao da saida ao aplicar pruning por magnitude e
quantizacao linear INT8 sobre operacoes centrais de Transformers.

\section{Estrategia de validacao}

A pesquisa separa quatro niveis de evidencia:

\begin{enumerate}
  \item \textbf{conceitual}: a definicao e o recorte correspondem a literatura;
  \item \textbf{matematica}: casos pequenos reproduzem valores calculados;
  \item \textbf{algoritmica}: implementacoes independentes produzem resultados equivalentes;
  \item \textbf{hardware}: o caminho executado, o tipo numerico, o formato e o kernel sustentam um ganho fisico mensuravel.
\end{enumerate}

A aprovacao de um nivel nao promove automaticamente o resultado ao nivel
seguinte. Em particular, zeros logicos e tensores dequantizados nao bastam para
afirmar aceleracao em hardware.

\section{Estado comprovado}

\begin{longtable}{@{}p{0.25\linewidth}p{0.18\linewidth}p{0.49\linewidth}@{}}
\toprule
\textbf{Componente} & \textbf{Estado} & \textbf{Evidencia principal} \\
\midrule
\endfirsthead
\toprule
\textbf{Componente} & \textbf{Estado} & \textbf{Evidencia principal} \\
\midrule
\endhead
Matematica da atencao & \statusok & Casos manuais, NumPy independente e testes deterministas. \\
Arquitetura do recorte & \statusok & Auditoria do bloco simplificado, dependencia entre posicoes e controle negativo. \\
Equivalencia entre frameworks & \statusok & Nove comparacoes entre NumPy, PyTorch manual, PyTorch SDPA e Keras/TensorFlow. \\
Pipeline experimental & \statusok & Configuracao versionada, duas execucoes, hashes, manifesto, metadados, analise e graficos. \\
Diagnostico em CPU & \statusok & Smoke e duas versoes do ponto $L=64$, $D=128$, sem conclusao de hardware. \\
Coleta principal em CUDA & \statuspending & Requer a maquina com RTX 4070 Ti Super. \\
Comparacao pruning versus INT8 & \statuspending & Depende da coleta principal e da auditoria do caminho efetivamente executado. \\
\bottomrule
\end{longtable}

\begin{icstatus}[Conclusao permitida neste momento]{icBlue}
Os scripts implementam e testam operacoes centrais de um bloco Transformer
simplificado e a atencao local e numericamente equivalente as referencias
escolhidas. Ainda nao ha evidencia experimental na GPU alvo para concluir que
pruning ou quantizacao produzam ganho real de hardware.
\end{icstatus}
